\chapter[Correcting Errors]{Correction of Throwing and Construction Errors}
\label{sec:correct}
\newpage

\section{Basic procedures}
There are many factors influencing the flight of a boomerang. It is impossible to predict the desired effect exactly, but with experience, you will have a fair idea how a boomerang will fly before it is thrown.

The only real promise of success in addition to various rules is usually: try, try and try again.
Rusty Harding claims: 
\begin{quote}
``Give me a rule of how to construct a boomerang and I will show you the boomerang to which it does not apply''.
\end{quote}

\begin{center}
\colorbox{BOrange!30}{\parbox{30em}{
\colorbox{BOrange!30}{\parbox{28em}{
	\color{black}
	\vspace{1em}
	When correcting throwing and construction errors, you should proceed step by step:
\begin{enumerate}
\item First you should try different throws.\\
Read about throwing and use section \ref{sec:DebugThrowing} to troubleshoot your throw.
\item ``Small'' non destructive adjustment\\
Bending of the wing tips, brake flaps, weighting with coins, smoothing down or roughing up the finish etc.
See Chapter \ref{sec:BasicImprovement}.
\item ``Big'' modifications\\
Undercuts, carving out, installation of weights, holes and slits.
See Chapter \ref{sec:BasicImprovement}.
\item Make another similar boomerang\\
Start the profiling differently. i.e. start with less trailing edge cutting, and a blunter leading edge.
\item Test changes in a different sequence.
\item  Try a different shape or material. 
\end{enumerate}
}}}}
\end{center}

See Chapter \ref{sec:BasicImprovement} for construction changes and chapter \ref{sec:ConstructionDisk} for using the construction disk. The disk is most useful after the simple throwing and tuning corrections have been tried, because it helps you see how a slight change in shape, airfoil, bend, weight or carving can move the average point of lift.
	
\section{How to achieve certain features}
\label{sec:AchieveCorrect}
\subsection{How to increase the effective lift}
\label{sec:increaseLift}
\begin{enumerate}
	\item Apply or increase the positive dihedral or angle of attack
	\item Carve out the underside (hollow)
	\item Undercut the leading edge
	\item Blunt leading edge 
	\item Make a longer trailing edge 
	\item Smooth the topside, make the underside rougher
	\item Make another boomerang with:
	\begin{itemize}
		\item Thicker material 
		\item Lighter material (to increase effective lift)
		\item More  distance  between the wing tips and the centre of rotation
	\end{itemize}
\end{enumerate}

\subsection{How to decrease the effective lift}
\label{sec:decreaseLift}
\begin{enumerate}
\item Apply or increase the negative dihedral or angle of attack 
\item Make the leading edge sharper
\item An undercut trailing edge
\item Make holes in the wing
\item Add flaps
\item Roughen the topside, smooth the underside
\item Make another boomerang with:
	\begin{itemize}
	\item Thinner material
	\item Shorter trailing edge 
	\item Heavier material
	\item Less distance between the wing tips and centre of rotation
	\item Narrower wing (reduced chord)
	\end{itemize}
\end{enumerate}

\subsection{How to increase rotation}
\begin{enumerate}
\item Throw with more spin 
\item Install weights near the wing tips
\item Reduce positive angle of attack or make it negative
\item Smooth the surfaces
\item Streamline the airfoils
\item Make another boomerang with:
	\begin{itemize}
	\item Thinner material
	\item Wider wings at the tips (increase chord)
\end{itemize}
\end{enumerate}

\subsection{How to decrease rotation}
\begin{enumerate}
\item Throw with less spin 
\item Apply or increase the positive angle of attack
\item Make the wing tips more blunt
\item Make the leading edge more blunt
\item Holes or slots in the wing especially near the tips
\item Add flaps near the wing tips
\item Make another boomerang with thicker material 
\end{enumerate}

\subsection{How to increase the range}
\begin{enumerate}
\item Throw with less tilt and harder
\item Add more weight near the wing tips
\item Decrease the effective lift (see section \ref{sec:decreaseLift}) 
\item Move the centre of rotation towards the wing tips (e.g. with weights on the tips)
\end{enumerate}

\subsection{How to reduce the range}
\begin{enumerate}
\item Throw with more tilt and less power
\item Remove weights
\item Increase the effective lift (see section \ref{sec:increaseLift}) 
\item Move the centre of rotation towards the elbow (e.g. with weights on the elbow)
\end{enumerate}

\subsection{How to increase the flight height}
\begin{enumerate}
\item Throw with more tilt 
\item Throw higher and harder 
\item Add more lift on arm 1 (with two-bladers) 
\item Apply or increase the positive dihedral
\item Add weights
\item Add flaps on the underside
\end{enumerate}

\subsection{How to reduce the flight height}
\begin{enumerate}
\item Throw with less tilt 
\item More lift on arm 2 (for two-bladers) 
\item Negative dihedral on both arms, to reduce ``diving'' add a slight positive angle of attack on both arms
\item Add flaps on the topside especially towards the trailing edge
\item Add weights
\item Drill a hole in the area of centre of rotation of a multiblader
\end{enumerate}

\subsection{How to increase the lay-over}
\begin{enumerate}
\item Increase dihedral or make it more positive
\item For two-bladers, increase the lift on arm 1 and decrease it on arm 2 
\item Increase the lift at the elbow; direction of air flow from the elbow towards the centre of rotation.
\end{enumerate}


\subsection{How to decrease the lay-over}
\begin{enumerate}
\item Reduce dihedral or make it more negative
\item For two-bladers, decrease the lift on Arm 1 and increase it on Arm 2.
\item Decrease the lift at the elbow; with the outside of the elbow as the leading edge
\end{enumerate}


